\subsection{Positive passages beyond a coefficient scale}
\label{sec:direct-adaptive-positive}

The bounded-radius estimates do not control positive graphs that leave
every fixed multiple of the coefficient scale. To treat these graphs
in the primary and secondary regions, we follow them within a small
unscaled radius box. An adapted height coordinate gives exact action
transport and estimates for the moving exit radius. Continuation from
that exit to a cone then permits the weighted-tail estimate to give
the positive endpoint slope. The corollary below compares it with the
negative endpoint slope to prove contraction after a sufficiently
large exit.

Write \(\varrho=\eta>0\) for the positive radius; in the negative
quadratic-drift application below, \(\varrho=s\rho\). Suppose that
\begin{equation}
\label{eq:direct-adaptive-drift}
 -q(\varrho)=\varrho^2(s+\varrho)\mathcal K(u,v,\varrho,p),\qquad
 u=\varrho_0/\varrho,\quad v=\varrho/(s+\varrho),\qquad s,\varrho_0\ge0,
\end{equation}
where \(\mathcal K\) has fixed positive upper and lower bounds
and uniform derivative bounds on an enlarged compact box of
the displayed base variables. Both \(s\) and the unscaled radius
are small. The coefficient \(\varrho_0\) is held fixed along a graph.
A restriction of the box in \(v\) is handled by a smooth
extension preserving positivity.
The factor \(\varrho^2(s+\varrho)=s\varrho^2+\varrho^3\) joins the
quadratic range \(\varrho\ll s\) to the cubic range
\(\varrho\gg s\).

Put
\[
 \varrho^+=s+\varrho,\qquad y=\frac{z}{\varrho^2\varrho^+},\qquad \mathfrak l=\log\varrho,
 \qquad \frac{d\tau}{d\mathfrak l}=\frac1{\varrho^+},\qquad
 \mathcal B=B(\varrho,\varrho^2\varrho^+y,p).
\]
The graph equation gives the extended system
\begin{equation}
\label{eq:direct-adaptive-system}
 \begin{aligned}
 u_\tau&=-(s+\varrho)u,& \varrho_\tau&=(s+\varrho)\varrho,& v_\tau&=sv,\\
 y_\tau&=\mathfrak a\,y+y^2\mathcal N,&
 \mathfrak a&=\mathcal K^{-1}-2s-3\varrho,&
 \mathcal N&=-\frac{\mathcal B}{\mathcal K(\mathcal K+\mathcal B y)} .
 \end{aligned}
\end{equation}
The defining relations among the base variables are invariant.
For a small fixed \(y_0>0\), the quantities
\(\mathfrak a,\ y_\tau/y,\ \partial_y y_\tau\) have fixed
positive lower bounds on \(|y|\le2y_0\).

\begin{lemma}[Adaptive height and cone landing]
\label{lem:direct-adaptive-positive}
For every fixed finite order \(k\), after reducing \(s\) and the
radius box, there is a uniformly regular coordinate
\[
 \Psi(u,v,\varrho,y,p)=y\exp\Theta(u,v,\varrho,y,p),\qquad
 \Theta=y\widetilde\Theta,\quad
 \|\widetilde\Theta\|_{C^k}\le C_k,
\]
with positive height derivative and bounded inverse derivatives,
such that \(\Psi_\tau=\mathfrak a\Psi\).
The adapted unscaled height \(\widehat z=\varrho^2(s+\varrho)\Psi\)
therefore satisfies exactly
\begin{equation}
\label{eq:direct-adaptive-action}
 \widehat z(\varrho)=
 \exp\left(\int_{\varrho_i}^{\varrho}\frac{t\,dt}{-q(t)}\right)
 \widehat z(\varrho_i)
\end{equation}
as long as the passage is thin, \(y\le y_0\).

For an exit at \(y=y_0\), let \(\varrho_e\) be its radius and
\(\lambda_{\rm in}=\log\widehat z(\varrho_i)\). Then
\begin{equation}
\label{eq:direct-adaptive-exit}
 -C(s+\varrho_e)\le\frac{d\log\varrho_e}{d\lambda_{\rm in}}
                 \le-c(s+\varrho_e),\qquad
 \left|\frac{d^j\log\varrho_e}{d\lambda_{\rm in}^j}\right|
                 \le C_j(s+\varrho_e)^j,\quad 1\le j\le k .
\end{equation}
There is a subsequent first landing on \(z=m_{\rm land}\varrho^2\), for a
small fixed \(m_{\rm land}>0\), with radius \(\varrho_l\) satisfying
\begin{equation}
\label{eq:direct-adaptive-landing}
 \begin{gathered}
 \varrho_e\le\varrho_l\le2\varrho_e,\qquad
 \partial_{\log\varrho_e}\log\varrho_l=1+O(\epsilon_e)>c>0,\\
 \bigl|\partial_{\log\varrho_e}^j\log\varrho_l\bigr|\le C_j\epsilon_e
       \quad(2\le j\le k),\qquad
 \epsilon_e=m_{\rm land}+(s+\varrho_e)|\log(s+\varrho_e)|,
 \end{gathered}
\end{equation}
with constants independent of \(m_{\rm land}\).
If a graph has \(y\ge y_0\) and \(z<m_{\rm land}\varrho^2\) at a radius
\(\bar\varrho\) of the box, its first landing after \(\bar\varrho\) has
radius at most \(2\bar\varrho\).
The construction applies to admissible graphs before their
endpoint.
\end{lemma}

\begin{proof}
\emph{The adapted height coordinate.}
At backward time \(t\ge0\), the base flow of
\eqref{eq:direct-adaptive-system} is explicitly
\[
 g_s(t)=\frac{1-\e^{-st}}s,\quad
 \varrho_b(t)=\frac{\varrho \e^{-st}}{1+\varrho g_s(t)},\quad
 u_b(t)=u\e^{st}(1+\varrho g_s(t)),\quad v_b(t)=v\e^{-st},
\]
where \(g_0(t)=t\). Since \(0\le g_s(t)\le t\) and the
denominator is at least one, each fixed-order state derivative
is bounded by a polynomial in \(t\) times \(\e^{jst}\).
The bounds are uniform for initial base points in the original compact
box. To cover backward trajectories that leave it, extend the
coefficients by smooth bounded retractions of the base variables into
the enlarged positive box, preserving the positive height-rate and
height-derivative margins. The extension agrees with the original
coefficients and makes \(\mathfrak a,\mathcal N\), together with their
required derivatives, bounded on the entire nonnegative base domain
and a small negative-height strip. Along the backward trajectories,
\(\varrho\) decreases.

Choose \(c_y>0\) so that \(\partial_y y_\tau\ge2c_y\)
on the extended box, and reduce \(s\) so that
\((k+1)s<c_y\). The base-flow bounds above verify
\cref{lem:direct-contracting-height-coordinate} with \(c_x=s\).
It supplies \(\Theta\) with the asserted bounds on
\(\Theta/y\) and the inverse, and \(\Psi_\tau=\mathfrak a\Psi\).
The logarithmic derivative of \(\varrho^2(s+\varrho)\) in \(\tau\) is
\(2s+3\varrho\), so
\((\log\widehat z)_{\mathfrak l}=1/(\varrho^+\mathcal K)\).
This is \eqref{eq:direct-adaptive-action}, by
\eqref{eq:direct-adaptive-drift}.

\emph{The exit radius.}
At the thin exit, \(\mathfrak l_e=\log\varrho_e\) solves
\[
 \mathcal M(\mathfrak l_e)=\lambda_{\rm in},\qquad
 \mathcal M(\mathfrak l)=
 \log[\varrho^2(s+\varrho)\Psi(u,v,\varrho,y_0,p)]
          -\int_{\varrho_i}^{\varrho}\frac{t\,dt}{-q(t)} .
\]
The transport equation gives
\[
 \mathcal M'(\mathfrak l)
 =-\frac{f_0}{\varrho^+}\,\partial_y\log\Psi(u,v,\varrho,y_0,p)
 =-\frac{g(u,v,\varrho,p)}{\varrho^+},\qquad c\le g\le C,
\]
where \(f_0=\mathfrak a y_0+y_0^2\mathcal N(u,v,\varrho,y_0,p)\) is the
height velocity at the exit. Fixed logarithmic radial derivatives of \(g\) are
bounded. The corresponding base vector field is
\(-u\partial_u+v(1-v)\partial_v+\varrho\partial_\varrho\), whose
fixed iterates applied to \(\varrho^+\) are \(O(\varrho^+)\). Set
\(\phi_0=-\varrho^+/g\), so that \(\partial_{\mathfrak l}^j\phi_0=O_j(\varrho^+)\).
Inverse differentiation gives
\[
 \frac{d^j\mathfrak l_e}{d\lambda_{\rm in}^j}
 =\bigl[(\phi_0\partial_{\mathfrak l})^{j-1}\phi_0\bigr]_{\mathfrak l=\mathfrak l_e}.
\]
Each expanded term contains exactly \(j\) factors, each a
radial derivative of \(\phi_0\) of some nonnegative order.
It is therefore \(O_j((s+\varrho_e)^j)\). The sign and two-sided
bound at first order follow from \(c\le g\le C\), proving
\eqref{eq:direct-adaptive-exit}.
Since the thin height increases and scalar order holds, the
exit is unique and its existence set is one
interval of incoming heights.

\emph{Passage to the cone.}
After exit use \(\mathcal Y=z/\varrho^2=\varrho^+y\) as the independent height.
The equations become
\[
 \mathcal Y_{\mathfrak l}=\frac{\mathcal Y}{B\mathcal Y+\varrho^+\mathcal K}-2\mathcal Y,\qquad
 \frac{d\mathfrak l}{d\mathcal Y}
 =\frac{\varrho^+\mathcal K+B\mathcal Y}
        {\mathcal Y(1-2\varrho^+\mathcal K-2B\mathcal Y)} .
\]
Choose \(m_{\rm land}\) small and then the radius box and \(s\) small,
so that the last denominator factor is at least \(1/2\)
on \(\mathcal Y\le m_{\rm land}\). Starting at \(\mathcal Y_e=(s+\varrho_e)y_0\), bootstrap
\(\varrho\le2\varrho_e\). With \(\varrho_e^+=s+\varrho_e\), integration gives
\[
 0\le\log(\varrho_l/\varrho_e)
 \le C\{m_{\rm land}+\varrho_e^+\log(m_{\rm land}/(y_0\varrho_e^+))\}<\log2 .
\]
This closes the bootstrap and proves arrival at \(\mathcal Y=m_{\rm land}\).
For the derivatives, write \(x=\log\varrho_e\) and
\(\mathfrak l(\mathcal Y;x)=\log\varrho(\mathcal Y;x)\). Denote by \(R(\mathfrak l,\mathcal Y)\) the
right side of the equation for \(\mathfrak l\) in the variable \(\mathcal Y\), with
parameters fixed, so that
\[
 \mathfrak l_{\mathcal Y}=R(\mathfrak l,\mathcal Y),\qquad \mathfrak l(\mathcal Y_e(x);x)=x.
\]
On the preceding bootstrap range,
\[
 |\partial_{\mathfrak l}^jR|\le C_j(\varrho_e^+/\mathcal Y+1),\qquad
 \int_{\mathcal Y_e}^{m_{\rm land}}(\varrho_e^+/\mathcal Y+1)\,d\mathcal Y
       \le C\epsilon_e,\quad
 \epsilon_e=m_{\rm land}+\varrho_e^+|\log \varrho_e^+|.
\]
Since the last denominator factor is at least \(1/2\), the constants
\(C_j\) depend only on the bounds for \(\mathcal K\) and \(B\), and not on
\(m_{\rm land}\).
The moving initial height is treated as in the moving-entrance part
of \cref{lem:direct-weighted-tail}: near the entrance set
\(\mathcal Y=\mathcal Y_e(x)\zeta\) and
\(\mathfrak l=x+\tilde{\mathfrak l}\). The rescaled equation has
right side \(\mathcal Y_eR\), whose required mixed derivatives are
\(O_j(\varrho_e^+)\) on a fixed \(\zeta\)-interval. Conversion back to
fixed \(\mathcal Y\) uses
\(\partial_x|_{\mathcal Y}=\partial_x|_\zeta-v_e\zeta\partial_\zeta\),
where \(v_e=\varrho_e/\varrho_e^+\) has bounded derivatives. Thus the
initial data for \(P_j=\partial_x^j\mathfrak l|_{\mathcal Y}\) are
\[
 P_1(\mathcal Y_e)=1-R(x,\mathcal Y_e)y_0\varrho_e
       =1+O(\varrho_e^+),\qquad
 P_j(\mathcal Y_e)=O_j(\varrho_e^+)\quad(j\ge2).
\]
Their triangular variational equations have linear coefficient
\(R_{\mathfrak l}\) and higher forcing coefficients
\(\partial_{\mathfrak l}^mR\). The displayed integral bound and
variation of constants therefore give
\(P_1(m_{\rm land})=1+O(\epsilon_e)\) and
\(P_j(m_{\rm land})=O_j(\epsilon_e)\) for \(j\ge2\), proving
\eqref{eq:direct-adaptive-landing}.

For an input above the thin threshold at a radius \(\bar\varrho\), the
same equation in \(\mathcal Y\), started at a larger value of \(\mathcal Y\), gives radial
growth to \(\mathcal Y=m_{\rm land}\) no larger than for an input at that
threshold; hence its landing radius is at most \(2\bar\varrho\).
An input with \(\mathcal Y\ge m_{\rm land}\) already has the required cone height.

\emph{Reaching the cone before the endpoint.}
The organizing endpoints satisfy
\(\eta_\pm^0=\pm a_*r_*+r_*^2\), in the notation of
\cref{lem:exterior-multiplier-new}, by
\eqref{eq:normal-rate-coordinate-new}. Hence
\(z_*/(\eta_\pm^0)^2\to1/2\) as \(z_*\downarrow0\).
For each sufficiently small fixed \(z_*\), the endpoint windows and
parameter neighborhood can therefore be chosen so that
\(\mathcal Y_*=z_*/\varrho_*^2\ge c_{\rm end}>0\), with one
\(c_{\rm end}\) independent of \(z_*\). Thus
\(m_{\rm land}<c_{\rm end}/2\) can be fixed before \(z_*\). The endpoint's normalized thin height is
\(y_*=\mathcal Y_* /(s+\varrho_*)>y_0\) on the small box in the coordinates
\((\eta,z)\).
Thus every admissible graph that starts thin exits and lands before
the fixed endpoint section. Monotonicity of the positive
radius keeps the construction in the chosen radius box.
\end{proof}

The preceding lemma controls the moving exit and its first cone
landing. To obtain a return-map estimate, we now differentiate the
endpoint reached from this landing. At
\(t=m_{\rm land}\varrho_l^2\), the entrance energy is \(U=t/(2m_{\rm land})\).
By the choice \(1/(2m_{\rm land})>4\sup B\), the moving-cone
part of \cref{lem:direct-weighted-tail} applies whenever
\eqref{eq:direct-tail-smallness} holds there. The endpoint energy
then has
\[
 \frac{dU(z_*;t/(2m_{\rm land}))}{dt}
       =(z_*/t)^{-2q_2}a(t),\qquad c<a(t)<C.
\]
Under this tail condition, suppose also that the adapted
incoming height satisfies
\(0<c\le d\lambda_{\rm in}/d\log h\le C\), where
\(\lambda_{\rm in}=\log\widehat z(\varrho_i)\). Then
\begin{equation}
\label{eq:direct-adaptive-source-slope}
 hI_{+,h}\asymp \varrho_e^2(s+\varrho_e)
                    (z_*/(m_{\rm land}\varrho_l^2))^{-2q_2}.
\end{equation}
Indeed \eqref{eq:direct-adaptive-exit} and
\eqref{eq:direct-adaptive-landing} give
\(-dt/d\log h\asymp t(s+\varrho_e)\), and \(I_+=z_*-U(z_*)\).

\begin{corollary}[Contraction after a large exit]
\label{cor:large-exit-contraction}
In the setting of \cref{lem:direct-adaptive-positive}, let
\(\mathcal J\subset\cH_p\) be an interval of central heights whose positive
half-orbits are thin at \(\varrho_i\), satisfy
\eqref{eq:direct-tail-smallness} at their landings, and have
\(c\le d\lambda_{\rm in}/d\log h\le C\). Let \(c_1\) be the constant of the
lower bound in \eqref{eq:direct-adaptive-source-slope}. Suppose that
\(q_2\le C_2s\) and \(0<m_-\le C_1s^3\) on \(\mathcal J\), with fixed
\(C_1,C_2\), and fix \(M_1\ge1\) with
\(c_1M_1^2(1+M_1)>2\Lambda_+C_1\). Then, for sufficiently small \(s\),
the central heights in \(\mathcal J\) with \(\varrho_e\ge M_1s\) form an
interval on which \(\mathcal L<0\). In particular \(\Delta_p\) is strictly
decreasing there and has at most one zero.
\end{corollary}

\begin{proof}
By \eqref{eq:direct-adaptive-exit} and the lower bound for
\(d\lambda_{\rm in}/d\log h\), the exit radius decreases with \(h\), so
these central heights form an interval. The landing precedes the
endpoint, so \(t=m_{\rm land}\varrho_l^2<z_*\), and the multiplier in
\eqref{eq:direct-adaptive-source-slope} is at least one if \(q_2\le0\).
If \(0<q_2\le C_2s\), then \(\varrho_l\ge\varrho_e\ge s\) gives
\(t\ge m_{\rm land}s^2\), and
\[
 (z_*/t)^{-2q_2}\ge\exp[-2C_2s\log(z_*/(m_{\rm land}s^2))]\ge\tfrac12
\]
for small \(s\). In both cases
\(m_+\ge\frac12c_1\varrho_e^2(s+\varrho_e)\), and therefore
\[
 \frac{m_-}{m_+}\le\frac{2C_1}{c_1}\,
        \frac1{(\varrho_e/s)^2(1+\varrho_e/s)} .
\]
The endpoints of central heights in \(\cH_p\) lie in the windows, where
\(\Pi_p'\le\Lambda_+\) by \cref{lem:exterior-multiplier-new}. For
\(\varrho_e\ge M_1s\), the choice of \(M_1\) gives
\(\e^{\mathcal L}=\Pi_p'(I_-)m_-/m_+<1\), and
\(\Delta_p'=I_+'(\e^{\mathcal L}-1)<0\) by
\eqref{eq:direct-completed-log-ratio}.
\end{proof}
